\documentclass[10pt]{article}
\usepackage[T1]{fontenc}
\usepackage[utf8]{inputenc}
\usepackage[french]{babel}
\usepackage{amsmath,amssymb,mathrsfs,mathpazo}
\usepackage{geometry}
\geometry{a4paper,margin=2cm}
\usepackage{graphicx,array,tabularx,tabularray,colortbl,multirow,multicol,booktabs,enumitem,pifont,ulem,fancybox,diagbox,cancel}
\usepackage{pstricks,pstricks-add,pst-plot,pst-tree,pst-node,pst-text,pst-func,pst-math,pst-bspline,pst-3dplot,pst-eucl,pst-solides3d}
\let\pscalculate\undefined
\usepackage{pst-bezier}
\usepackage{tikz,pgfplots,tkz-tab}
\usetikzlibrary{arrows,arrows.meta,patterns,calc,positioning,trees,decorations.markings,shapes.geometric,shapes.misc,shapes.arrows,intersections,decorations.pathmorphing,decorations.pathreplacing,backgrounds,shadows,babel}
\pgfplotsset{compat=1.18}
\usepgfplotslibrary{fillbetween}
\usepackage[np]{numprint}
\usepackage{listings,makecell,pas-tableur,twoopt}
\let\barre\undefined
\usepackage[pstricks]{tablvar}
\definecolor{darkgreen}{RGB}{0,100,0}
\providecommand{\red}{\color{red}}
\providecommand{\blue}{\color{blue}}
\providecommand{\green}{\color{green}}
\providecommand{\black}{\color{black}}
\providecommand{\white}{\color{white}}
\providecommand{\gray}{\color{gray}}
\providecommand{\cyan}{\color{cyan}}
\providecommand{\magenta}{\color{magenta}}
\providecommand{\yellow}{\color{yellow}}
\providecommand{\violet}{\color{violet}}
\providecommand{\orange}{\color{orange}}
\providecommand{\euro}{\texteuro}
\providecommand{\EUR}{\texteuro}
\providecommand{\ieme}{\textsuperscript{e}}
\providecommand{\degres}{\ensuremath{^\circ}}
\providecommand{\degree}{\ensuremath{^\circ}}
\providecommand{\vv}[1]{\overrightarrow{#1}}
\providecommand{\Coord}[1]{#1}
\providecommand{\decosix}{}
\providecommand{\starredbullet}{\ensuremath{\bullet}}
\pagestyle{empty}
\setlength{\parindent}{0pt}
\usepackage{adjustbox}\DeclareUnicodeCharacter{2236}{:}\DeclareUnicodeCharacter{2190}{\ensuremath{\leftarrow}}
\providecommand{\R}{\mathbb{R}}
\providecommand{\N}{\mathbb{N}}
\providecommand{\Z}{\mathbb{Z}}
\providecommand{\Q}{\mathbb{Q}}
\providecommand{\D}{\mathbb{D}}
\providecommand{\C}{\mathbb{C}}
\providecommand{\e}{{\mathrm{e}}}
\providecommand{\vect}[1]{\overrightarrow{#1}}
\providecommand{\vectt}[1]{\overrightarrow{\mathrm{#1}}}
\providecommand{\vvt}[1]{\overrightarrow{\mathrm{#1}}}
\providecommand{\np}[2][]{{#2~#1}}
\providecommand{\vv}[1]{{\overrightarrow{#1}}}
\renewcommand{\d}{{\mathrm{d}}}
\providecommand{\degres}{{^\circ}}
\providecommand{\degre}{{^\circ}}
\providecommand{\blue}{\color{blue}}
\providecommand{\red}{\color{red}}
\providecommand{\enskip}{\quad}
\providecommand{\llbracket}{\lbrack\!\lbrack}
\providecommand{\rrbracket}{\rbrack\!\rbrack}
\providecommand{\euro}{{\text{€}}}
\providecommand{\no}{{\text{n°}}}
\providecommand{\interval}[2]{{[#1;#2]}}
\providecommand{\intervFO}[2]{{[#1;#2[}}
\providecommand{\dotfill}{\ldots}

\providecommand{\ds}{\displaystyle}
\providecommand{\textding}[1]{#1}
\providecommand{\Oij}{\(\mathrm{O};\vec{\imath},\vec{\jmath}\)}
\providecommand{\Oijk}{\(\mathrm{O};\vec{\imath},\vec{\jmath},\vec{k}\)}
\providecommand{\Ouv}{\(\mathrm{O};\vec{u},\vec{v}\)}
\providecommand{\intervFF}[2]{\ensuremath{[#1;#2]}}
\renewcommand{\euro}{\ifmmode\text{\texteuro}\else\texteuro\fi}
\DeclareUnicodeCharacter{00B0}{\ifmmode^{\circ}\else\textdegree\fi}
\DeclareUnicodeCharacter{20AC}{\ifmmode\text{\texteuro}\else\texteuro\fi}
% Énoncés : Éducation nationale. Transcription et corrigés : APMEP, Denis Vergès et contributeurs. Découpage et classement : CoopMaths, Sébastien Lozano, Éric Elter et contributeurs. Adaptation de la présentation : Maths Explorer.
% CC BY-SA : https://coopmaths.fr/www/about/
% https://coopmaths.fr/alea/static/bac/2025/tex/bac_2025_06_sujet1_asie_4_cor.tex
\begin{document}
\noindent\textit{APMEP / CoopMaths. CC BY-SA. Présentation adaptée par Maths Explorer.}\par\medskip
\medskip

On considère $f$ la fonction définie sur l'intervalle $]0; +\infty[$ par $f(x) = \dfrac{\e^{\sqrt{x}}}{2 \sqrt{x}}$ et on appelle $C_f$ sa courbe représentative dans un repère orthonormé.

\begin{enumerate}
\item On définit la fonction $g$ sur l'intervalle $]0~;~ +\infty[$ par $g(x) = \e^{\sqrt{x}}$
	\begin{enumerate}
		\item La fonction $g$ est continue et dérivable sur $]0~;~+\infty[$, donc $\forall x \in ]0~;~+\infty[$,

$g'(x) = \dfrac{1}{2\sqrt{x}}\times \e^{\sqrt{x}}= \dfrac{\e^{\sqrt{x}}}{2\sqrt{x}} = f(x)$.
		\item La fonction $f$ est continue et dérivable sur $]0~;~+\infty[$, donc $\forall x \in ]0~;~+\infty[$,
	
$f'(x) = \dfrac{\dfrac{1}{2\sqrt{x}}\e^{\sqrt{x}}\times2\sqrt{x}-\e^{\sqrt{x}}\times 2\dfrac{1}{2\sqrt{x}}}{(2\sqrt{x})^2} = \dfrac{\e^{\sqrt{x}}-\dfrac{\e^{\sqrt{x}}}{\sqrt{x}}}{4x}=\dfrac{\dfrac{\sqrt{x}\e^{\sqrt{x}}-\e^{\sqrt{x}}}{\sqrt{x}}}{4x}= \dfrac{\e^{\sqrt{x}} (\sqrt{x} - 1)}{4 x \sqrt{x}}$.
	\end{enumerate}
	\item 
	\begin{enumerate}
		\item $\ds \lim\limits_{x \to 0} \e^{\sqrt{x}}=1$ et $\ds \lim\limits_{x \to 0} 2\sqrt{x}=0$ avec $\dfrac{\e^{\sqrt{x}}}{2 \sqrt{x}} > 0$ donc $\ds \lim\limits_{x \to 0}f(x)=+\infty$.
		\item Donc $C_f$ admet une asymptote verticale d'équation $x=0$.
	\end{enumerate}
	\item 
	\begin{enumerate}
		\item D'après les croissances comparées, $\ds \lim\limits_{u \to +\infty} \dfrac{\e^u}{u}=+\infty$ donc par composition de limites (en posant $u(x)=\sqrt{x}$) $\lim\limits_{x \to +\infty} f(x)=+\infty$.
		\item $\forall x \in ]0~;~+\infty[$, $\e^{\sqrt{x}}>0$ et $4 x \sqrt{x}>0$ donc $f'(x)$ a le même signe que $\sqrt{x}-1$.

$\sqrt{x}-1\geqslant 0 \iff \sqrt{x}\geqslant1 \iff x\geqslant1$

\begin{center}
	\begin{tikzpicture}
			\tkzTab{$x$ / 1 , $f'(x)$ / 1, $f(x)$ / 1.5}{$0$, $1$, $+\infty$}
			{d, -, z, +, }
			{D+ / $+\infty$, -/ $\dfrac{\e}{2}$, +/ $+\infty$}
	\end{tikzpicture}
\end{center}
		\item La fonction $f$ est continue et strictement croissante sur l'intervalle $[1~;~+\infty[$ à valeurs dans $\left[\dfrac{\e}{2}~;~+\infty\right[$. Or $2\in \left[\dfrac{\e}{2}~;~+\infty\right[$ (car $\dfrac{\e}{2}<2$) donc d'après le corollaire du TVI, l'équation $f(x) = 2$ admet une unique solution notée $\alpha$ dans $[1;+\infty[$.

À la calculatrice : $\alpha\approx 4,6$.
	\end{enumerate}
	\item $I = \ds\int_1^2 f(x) \, \text{d}x$.
	\begin{enumerate}
		\item $\forall x \in [1~;~2]$, $f(x)=\dfrac{\e^{\sqrt{x}}}{2 \sqrt{x}}=\dfrac{1}{2\sqrt{x}} \times {\e^{\sqrt{x}}}$.
En remarquant que la dérivée de la fonction $g$ est la fonction $f$, $g$ est une primitive de $f$.

Donc $I = \ds\int_1^2 f(x) \, dx=\left[g(x)\right]_1^2=\left[\e^{\sqrt{x}}\right]_1^2= \e^{\sqrt{2}}- \e$
		\item L'aire du domaine $\begin{cases} 1\leqslant x \leqslant 2 \\ 0\leqslant y \leqslant f(x)\end{cases}$ est égale à $\e^{\sqrt{2}}- \e$.
		
$I$ est donc l'aire du domaine compris entre $\mathcal{C}_f$ , l'axe des abscisses et les droites d'équation $x =1$ et $x =2$.
	\end{enumerate}
	\item $\forall x \in ]0~;~ +\infty[$, $f''(x) = \dfrac{\e^{\sqrt{x}} (x - 3 \sqrt{x} + 3)}{8 x^2 \sqrt{x}}$
	\begin{enumerate}
		\item On pose $\forall x \in ]0~;~ +\infty[$, $X=\sqrt{x}$.

$x - 3 \sqrt{x} + 3=X^2-3X+3$.

		$\Delta=(-3)^2-4\times 1\times 3=-3<0$ donc $\forall X \in ]0~;~ +\infty[$, $X^2-3X+3>0$ donc $\forall x \in ]0; +\infty[$, $x - 3 \sqrt{x} + 3>0$.
		\item $\forall x \in ]0; +\infty[$, $e^{\sqrt{x}}>0$ et $x - 3 \sqrt{x} + 3>0$ et $8 x^2 \sqrt{x}>0$ donc $f''(x)>0$.
		
		La fonction $f$ est convexe sur $]0~;~ +\infty[$.
	\end{enumerate}
\end{enumerate}

\end{document}
